\subsection{在体论中的应用}

命题 9. —— 设 L 是一个体，E 是 \(\operatorname{End}_{\mathbf{Z}}(\mathrm{L})\) 的一个子环，它含映射 \(\gamma_{a}: x \mapsto ax\)（对每个 \(a \in L\)）。用 K 表示 L 中满足下列条件的元素 a 组成的集合：对 L 中每个 x 及 E 中每个 u，\(u(xa) = u(x)a\)；它是 L 的一个子体。

设 V 是右 K-向量空间 L 的一个有限维线性子空间，h 是一个从 V 到 L 的 K-线性映射。存在 E 的一个元素在 V 上与 h 相合。

我们把 L 视为左 E-模。由于 E 含诸左乘 \(\gamma_{a}\)，L 的每个 E-子模都是体 L 的左理想；因此 E-模 L 是单模。加群 L 的与诸 \(\gamma_{a}\) 交换的每个自同态都形如 \(\delta_{b}: x \mapsto xb\)，其中 b 属于 L。于是 \(b \mapsto \delta_{b}\) 是从 K 到 \(\operatorname{End}_{\operatorname{E}}(\operatorname{L})\) 的反环的一个同构，而后者是一个体。

因此 E-模 L 的双交换子 \(E''\) 由右 K-向量空间 L 的自同态组成。设 v 是 K-向量空间 L 的一个自同态，它在 V 上的限制与 h 相合；它是 \(E''\) 的元素。设 \((x_{i})_{i\in\mathrm{I}}\) 是 V 在 K 上的一个基；由定理 3（VIII, 第 88 页），存在 E 的元素 u，使得 \(u(x_{i}) = v(x_{i}) = h(x_{i})\)（对每个 \(i \in I\)）。由线性性推得，对 V 中每个 x 有 \(u(x) = h(x)\)。

推论. —— 设 L 是一个体。设 \(\Gamma\) 是体 L 的自同构群的一个子群，K 是 \(\Gamma\) 的不变量体。设 V 是 L 的一个右 K-线性子空间，它在 K 上有有限维 n。则存在元素 \(\sigma_{1},\ldots,\sigma_{n}\)（属于 \(\Gamma\)），具有下列性质：对每个从 V 到 L 的 K-线性映射 u，存在 L 的元素 \(a_{1},\ldots,a_{n}\)，使得对 V 中每个 x 有 \(u(x)=\sum_{i=1}^{n}a_{i}\sigma_{i}(x)\)。

用 E 表示形如 \(x \mapsto \Sigma_{\sigma \in \Gamma} a_{\sigma} \sigma(x)\) 的从 L 到 L 的映射组成的集合，其中 \((a_{\sigma})_{\sigma \in \Gamma}\) 是 L 的元素组成的有限支集族。对 L 中每个 a 我们有 \(\gamma_{a} \in E\)，且 E 是加群 L 的自同态环的一个子环。此外，体 K 由 L 中满足下列条件的元素 a 组成：对 L 中每个 x 及 E 中每个 u，\(u(xa) = u(x)a\)。

设 H 是左 L-向量空间 \(\operatorname{Hom}_{\mathrm{K}}(\mathrm{V},\mathrm{L})\)；它具有维数 n。由命题 9，它由 \(\Gamma\) 的元素在 V 上的限制生成。存在 n 个元素 \(\sigma_{1},\ldots,\sigma_{n}\)（属于 \(\Gamma\)），它们在 V 上的限制构成 H 在 L 上的一个基。本推论由此得出。

注. —— 当体 L 交换时，本推论化为阿廷定理（V, §10, No.~6, 第 65 页）。

